% =============================================================================
% Capitulo 6 - Conclusiones y lineas futuras
% =============================================================================
\chapter{Conclusiones y líneas futuras}
\label{cap:conclusiones}

\section{Conclusiones}
\label{sec:conclusiones}

Este Trabajo Fin de Grado ha evolucionado desde un enfoque inicial de clasificación global de espectrogramas hacia una arquitectura más robusta de detección e identificación en dos etapas. La primera etapa, basada en YOLOv8, localiza ráfagas en el espectrograma y las clasifica como \emph{drone} o \emph{interference}. La segunda etapa, basada en clasificadores como ResNet18, se reserva para la identificación de modelos de dron una vez detectada la presencia de señal compatible con UAV.

La principal conclusión experimental es que la firma tiempo-frecuencia del DJI Mini 3 es separable morfológicamente de las interferencias consideradas en la banda de \SI{2.4}{\giga\hertz}. El detector YOLOv8n entrenado durante 60 épocas alcanza en test una precisión global de 0,760, un recall global de 0,391, un mAP@50 de 0,388 y un mAP@50-95 de 0,203. Aunque el recall es limitado, la matriz de confusión normalizada no muestra confusión cruzada entre las clases activas \emph{drone} e \emph{interference}. Esto significa que, cuando el modelo emite una detección positiva, distingue correctamente la morfología de ambas clases.

El cambio metodológico resulta relevante. El baseline de clasificación global podía aprender sesgos contextuales de la captura, mientras que YOLO obliga al modelo a aprender regiones locales de señal. Esto hace que el resultado sea más interpretable y más coherente con el fenómeno físico que se desea detectar: ráfagas FHSS/OcuSync en el dominio tiempo-frecuencia.

\section{Limitaciones}
\label{sec:limitaciones}

La primera limitación es el recall del detector. El modelo no detecta todos los bursts anotados, por lo que todavía no puede considerarse un contador exhaustivo de ráfagas. Esta limitación puede deberse al pequeño tamaño de muchas regiones, al redimensionamiento de la imagen a \num{640}x\num{640}, al uso de YOLOv8n como arquitectura ligera y al carácter automático de las etiquetas.

La segunda limitación es el etiquetado débil. Las cajas se generan mediante umbrales y filtros morfológicos, no mediante anotación manual. Aunque este procedimiento es coherente y reproducible, puede producir etiquetas incompletas o inconsistentes en casos frontera. Una validación manual de un subconjunto permitiría estimar mejor la calidad del \emph{ground truth}.

La tercera limitación es la heterogeneidad de la clase \emph{interference}. Esta clase agrupa Bluetooth, WiFi residual y ruido, por lo que resulta más difícil de modelar que la clase \emph{drone}, asociada a una firma más homogénea del DJI Mini 3.

\section{Líneas futuras}
\label{sec:lineas_futuras}

Como trabajo futuro inmediato se propone entrenar variantes de mayor capacidad, como YOLOv8s o YOLOv8m, para comprobar si aumenta el recall sin introducir confusión entre clases activas. También sería conveniente realizar una anotación manual parcial de espectrogramas para validar el auto-etiquetado y medir cuánto afecta el ruido de etiquetas al rendimiento.

Otra línea directa es completar la segunda etapa de identificación de modelos de dron. Para ello se pueden utilizar capturas de distintos UAV, generar espectrogramas o recortes a partir de las regiones detectadas, y entrenar un clasificador multiclase como ResNet18. Esta etapa permitiría pasar de un sistema de detección de presencia a un sistema de detección e identificación.

Finalmente, el sistema puede extenderse hacia una versión operativa de búsqueda y confirmación. En dicha versión, GNU Radio barrería frecuencias de interés, Python generaría espectrogramas cortos y YOLO decidiría si bloquear la frecuencia para realizar una captura más larga. Este esquema conectaría el detector offline desarrollado en el TFG con un demostrador práctico de vigilancia RF.
